\documentclass[conference]{IEEEtran}
\IEEEoverridecommandlockouts
% Template version as of 6/27/2024

\usepackage{cite}
\usepackage{amsmath,amssymb,amsfonts}
\usepackage{algorithmic}
\usepackage{graphicx}
\usepackage{textcomp}
\usepackage{xcolor}
\usepackage{booktabs}
\def\BibTeX{{\rm B\kern-.05em{\sc i\kern-.025em b}\kern-.08em
		T\kern-.1667em\lower.7ex\hbox{E}\kern-.125emX}}

\begin{document}

	\title{Securing Image Encryption via a Hybrid Chaos-PQC Scheme Utilizing ML-KEM and Multi-Chaos}

\author{
	\begin{tabular}{cc}
		\begin{tabular}{c}
			\IEEEauthorblockN{Hafas Furqan \\
				\textit{Cryptographic Engineering} \\
				\textit{Politeknik Siber dan Sandi Negara}\\
				Bogor, Indonesia \\
				hafas.furqan@student.poltekssn.ac.id}
		\end{tabular}
		&
		\begin{tabular}{c}
			\IEEEauthorblockN{Yeni Farida \\
				\textit{Cryptographic Engineering} \\
				\textit{Politeknik Siber dan Sandi Negara}\\
				Bogor, Indonesia \\
				yeni.farida@poltekssn.ac.id}
		\end{tabular}
		\\[1.5cm]
		\begin{tabular}{c}
			\IEEEauthorblockN{Andika Aryansyach Fauzan \\
				\textit{Cryptographic Engineering} \\
				\textit{Politeknik Siber dan Sandi Negara}\\
				Bogor, Indonesia \\
				andika.aryansyach@student.poltekssn.ac.id}
		\end{tabular}
		&
		\begin{tabular}{c}
			\IEEEauthorblockN{Aga Ucu Fradana \\
				\textit{Cryptographic Engineering} \\
				\textit{Politeknik Siber dan Sandi Negara}\\
				Bogor, Indonesia \\
				aga.ucu@student.poltekssn.ac.id}
		\end{tabular}
	\end{tabular}
}

	\maketitle

	\begin{abstract}
		Chaos-based image encryption schemes offer strong statistical security properties but remain vulnerable to quantum attacks when their key establishment relies on classical public-key primitives such as RSA or ECC. No prior scheme combines a NIST-standardized post-quantum key encapsulation mechanism (KEM), cryptographically independent chaos parameter derivation, and a multi-map confusion-diffusion pipeline into a unified architecture. This paper proposes and evaluates a hybrid Chaos-PQC image encryption scheme that addresses this gap. Key establishment is performed using ML-KEM-768 (FIPS~203, NIST security category~3), which produces a 32-byte session-unique shared secret secured by the Module Learning With Errors (MLWE) hardness assumption. The shared secret is expanded by HKDF-SHA-256 into 144 bytes of output keying material, partitioned into 18 cryptographically independent parameters that seed a four-stage confusion-diffusion-confusion-diffusion pipeline: the Lorenz Attractor provides two-dimensional pixel permutation at each confusion stage, the Logistic Map generates the first-stage diffusion keystream, and the Tent Map---selected for its near-uniform output distribution---generates the second-stage keystream. Experiments on nine grayscale PNG images across three resolution classes ($512\!\times\!512$, $900\!\times\!900$, $1080\!\times\!1080$) demonstrate mean information entropy of 7.9994~bits (deviation from ideal: 0.0006), mean absolute inter-pixel correlation of $7\!\times\!10^{-4}$, and mean PSNR of 7.79~dB, confirming strong and content-independent encryption quality. Bit-exact decryption was verified for all images; a single-bit parameter fault rendered decryption completely infeasible, validating chaotic sensitivity.
	\end{abstract}

	\begin{IEEEkeywords}
		image encryption, post-quantum cryptography, ML-KEM, FIPS~203, HKDF, Lorenz attractor, Logistic Map, Tent Map, confusion-diffusion, chaos-based cryptography
	\end{IEEEkeywords}

% -------------------------------------------------------
\section{Introduction}
% -------------------------------------------------------

The rapid growth of digital public services---radiology archiving systems, telemedicine platforms, and biometric management applications---has substantially intensified the exchange of sensitive visual data. A large-scale audit by Greenbone Networks reported that misconfigured PACS servers left over 1.2 billion medical images accessible without authentication from the public internet \cite{b1}. In Indonesia, the national health insurance agency (BPJS Kesehatan) suffered a breach exposing identity numbers, phone numbers, addresses, and e-mail accounts of millions of citizens, with data subsequently traded in underground forums \cite{b2,b3}; the Ministry of Social Affairs likewise reported a breach of 102 million personal records \cite{b3}. Indonesia's Law No.~27 of 2022 on Personal Data Protection designates health, biometric, genetic, and children's data as sensitive categories requiring strong protection mechanisms \cite{b4}.

Chaos-based image encryption is widely studied due to three cryptographically desirable properties: \textit{sensitivity to initial conditions} (a minuscule parameter change yields a completely different output), \textit{ergodicity} (the trajectory uniformly visits all regions of the state space), and \textit{pseudo-randomness} (the output is statistically indistinguishable from true randomness) \cite{b9,b10}. The \textit{confusion-diffusion} architecture \cite{b12} is the predominant structural framework: confusion scrambles pixel positions to destroy spatial correlation, while diffusion modifies pixel intensity values so that local changes propagate globally. An ideal cipher image exhibits a near-uniform intensity histogram, inter-pixel correlation approaching zero, and entropy near the theoretical maximum of 8~bits/symbol \cite{b15}.

Single-map schemes suffer inherent limitations: the Logistic Map achieves full chaos only for $r \in (3.57, 4]$ and produces a non-uniform keystream distribution that constrains the effective key space \cite{b25}. Lawnik formally demonstrated that combining the Logistic and Tent Maps produces a composite system with superior statistical properties \cite{b25}. Praveen et al.\ consolidated these insights into a dual confusion--dual diffusion architecture using the Lorenz Attractor, Logistic Map, and Tent Map, outperforming single-map schemes \cite{b26,b29}. However, all such works rely on classically exchanged parameters vulnerable to Shor's algorithm on a quantum computer \cite{b18}.

The security of a chaos-based scheme depends not only on the quality of the chaotic dynamics but critically on the mechanism by which seeding parameters are generated and managed \cite{b20}. NIST finalized ML-KEM (FIPS~203) in August 2024 \cite{b6}, grounded in the hardness of the MLWE problem, believed to resist both classical and quantum attacks. A single 32-byte ML-KEM shared secret must be expanded via HKDF \cite{b7,b8} to derive multiple independent chaos parameters without key reuse. He et al.\ \cite{b20} demonstrated HKDF integration with chaos encryption but used ECC (quantum-vulnerable); Shakib et al.\ \cite{b23} confirmed ML-KEM+HKDF feasibility without image encryption. No prior work unifies all three components---a NIST-standardized PQC KEM, HKDF-based parameter derivation, and a multi-map confusion-diffusion pipeline---into a single coherent hybrid scheme. This paper closes that gap.

The main contributions are:
\begin{enumerate}
	\item A \textit{hybrid Chaos-PQC image encryption scheme} integrating ML-KEM-768 (FIPS~203, security category~3) for post-quantum key establishment with a four-stage multi-map chaos encryption pipeline.
	\item An \textit{HKDF-SHA-256 parameter derivation layer} mapping the ML-KEM shared secret into four cryptographically independent, context-bound subkeys, ensuring session-unique, non-reusable seeds for each chaos component.
	\item A \textit{dual confusion--dual diffusion encryption architecture} combining the Lorenz Attractor (1st confusion), Logistic Map (1st diffusion), Lorenz Attractor with a distinct HKDF-derived trajectory (2nd confusion), and Tent Map (2nd diffusion), evaluated on nine grayscale images across three resolution classes.
	\item \textit{Statistical evaluation} demonstrating entropy ${>}\,7.998$ bits/symbol, absolute inter-pixel correlation ${<}\,0.004$, and PSNR ${<}\,9$~dB across all cipher images.
\end{enumerate}

% -------------------------------------------------------
\section{Background and Related Work}
% -------------------------------------------------------

\subsection{Confusion-Diffusion Paradigm}

A grayscale image is represented as an $H \times W$ pixel matrix, where each element stores an 8-bit intensity value in $[0, 255]$ \cite{b13}. Image encryption transforms plain image $P$ into a visually unintelligible cipher $E$ through a cryptographically keyed process \cite{b12}. Confusion scrambles the spatial arrangement of pixels---breaking inter-pixel correlation---while diffusion modifies intensity values through a keystream XOR operation, ensuring that a one-bit plaintext change propagates to the entire cipher image. Quality is assessed by histogram uniformity, inter-pixel correlation approaching zero \cite{b14}, and entropy $H = -\sum_{v=0}^{255} p(v)\log_2 p(v)$ approaching 8~bits/symbol \cite{b15}. Low PSNR (${<}\,10$~dB) confirms high visual dissimilarity \cite{b12}.

\subsection{ML-KEM (FIPS 203)}

Classical key exchange protocols (RSA, Diffie-Hellman) are broken in polynomial time by Shor's algorithm on a large-scale quantum computer \cite{b18}. ML-KEM (FIPS~203) \cite{b6}, derived from CRYSTALS-Kyber, provides quantum-resistant key encapsulation based on the MLWE problem. Three variants balance security and bandwidth (Table~\ref{tab:mlkem}); this work uses ML-KEM-768 (NIST category~3), producing a 32-byte shared secret $K$. The protocol operates as follows \cite{b6}: (1)~\texttt{ML-KEM.KeyGen()} generates encapsulation key $\mathit{ek}$ (public) and decapsulation key $\mathit{dk}$ (private); (2)~\texttt{ML-KEM.Encaps($\mathit{ek}$)} produces ciphertext $c$ and shared secret $K$; (3)~\texttt{ML-KEM.Decaps($\mathit{dk},c$)} deterministically recovers the identical $K$, with correctness guaranteed with overwhelming probability.

\begin{table}[htbp]
	\caption{ML-KEM Parameter Sets \cite{b5,b27}}
	\label{tab:mlkem}
	\begin{center}
		\begin{tabular}{|l|c|c|c|c|}
			\hline
			\textbf{Variant} & \textbf{NIST Cat.} & \textbf{ek (B)} & \textbf{dk (B)} & \textbf{ct (B)} \\
			\hline
			ML-KEM-512  & 1 & 800  & 1632 & 768  \\
			ML-KEM-768  & 3 & 1184 & 2400 & 1088 \\
			ML-KEM-1024 & 5 & 1568 & 3168 & 1568 \\
			\hline
		\end{tabular}
	\end{center}
\end{table}

\subsection{HMAC Key Derivation Function (HKDF)}

Directly using the 32-byte ML-KEM output as a seed for multiple chaotic maps constitutes key reuse, creating cross-component correlation vulnerabilities \cite{b19}. HKDF \cite{b8}, recommended by NIST SP~800-56C Rev.~2 \cite{b7}, resolves this through an Extract-then-Expand construction: (1)~\textit{Extract} computes $\mathrm{PRK} = \mathrm{HMAC}(\mathit{salt},\,\mathit{IKM})$, producing a pseudorandom key from the shared secret; (2)~\textit{Expand} generates chained blocks $T(i) = \mathrm{HMAC}(\mathrm{PRK},\,T(i{-}1)\,\|\,\mathit{info}\,\|\,i)$, concatenated to form the output keying material (OKM). Using distinct \texttt{info} labels per chaos component ensures the resulting subkeys are mutually independent and context-bound, preventing cross-use across sessions \cite{b19}.

\subsection{Chaotic Maps}

\subsubsection{Lorenz Attractor}
The Lorenz system, introduced in 1963 \cite{b28}, is a 3-D continuous-time nonlinear system:
\begin{equation}
	\dot{x} = \sigma(y-x),\quad \dot{y} = x(\rho-z)-y,\quad \dot{z} = xy-\beta z,
	\label{eq:lorenz_bg}
\end{equation}
whose aperiodic trajectories---sensitive to initial conditions (the Butterfly Effect)---produce high-quality pixel permutation indices when argsorted \cite{b29}.

\subsubsection{Logistic Map}
The Logistic Map is a 1-D discrete-time chaotic map \cite{b11}:
\begin{equation}
	x_{n+1} = r\,x_n(1-x_n), \quad x_0 \in (0,1),
	\label{eq:logistic_bg}
\end{equation}
exhibiting full chaos for $r \in (3.57, 4]$. Its output concentrates near the interval boundaries, producing a non-uniform keystream distribution that limits its direct use as a sole diffusion source and constrains the effective key space \cite{b25}.

\subsubsection{Tent Map}
The Tent Map is a 1-D piecewise-linear chaotic map \cite{b9}:
\begin{equation}
	x_{n+1} = \begin{cases} x_n/\mu, & x_n < \mu \\ (1-x_n)/(1-\mu), & x_n \geq \mu \end{cases}
	\label{eq:tent_bg}
\end{equation}
with $\mu \in (0,1)$. Unlike the Logistic Map, it produces a nearly uniform output distribution over $[0,1)$ \cite{b25}, making it the preferred choice for the final diffusion stage where a flat cipher-image histogram is required. Lawnik \cite{b25} formally demonstrated that the Tent Map complements the Logistic Map in a combined system.

\subsection{Related Work and Research Gap}

Table~\ref{tab:related} summarizes the most directly relevant prior works and their gaps.

\begin{table}[htbp]
	\caption{Summary of Related Work}
	\label{tab:related}
	\begin{center}
		\begin{tabular}{|p{1.4cm}|c|p{2.0cm}|p{2.0cm}|}
			\hline
			\textbf{Reference} & \textbf{Year} & \textbf{Key Mechanism} & \textbf{Gap} \\
			\hline
			He et al. \cite{b20} & 2025 & ECC + HKDF + chaos & ECC is quantum-vulnerable \\
			\hline
			Shakib et al. \cite{b23} & 2025 & ML-KEM + HKDF & No image encryption \\
			\hline
			Praveen et al. \cite{b26,b29} & 2023 & Lorenz + Logistic + Tent & No PQC keying \\
			\hline
			Qayyum et al. \cite{b22} & 2020 & Two chaotic maps & No PQC; single diffusion \\
			\hline
			Lawnik \cite{b25} & 2018 & Logistic + Tent analysis & No encryption scheme \\
			\hline
		\end{tabular}
	\end{center}
\end{table}

He et al.\ \cite{b20} demonstrated that HKDF produces mutually independent subkeys from a single shared secret, but their key establishment relies on ECC, remaining vulnerable to Shor's algorithm \cite{b18}. Shakib et al.\ \cite{b23} confirmed ML-KEM keys can be safely extended via HKDF in practice, but do not address image encryption. Praveen et al.\ \cite{b26,b29} established the dual confusion--dual diffusion pipeline as effective for grayscale images, achieving superior statistical properties over single-map schemes; however, chaos parameters are user-provided, placing key establishment outside a post-quantum security boundary. Qayyum et al.\ \cite{b22} and Lawnik \cite{b25} further confirmed the superiority of multi-map approaches without addressing quantum-resistant key management.

The collective gap is clear: no existing scheme combines (1) a NIST-standardized PQC KEM, (2) HKDF-based derivation of independent chaos parameters, and (3) a multi-map confusion-diffusion pipeline into a unified hybrid architecture. The present work addresses this gap.

% -------------------------------------------------------
\section{Proposed Hybrid Chaos-PQC Scheme}
% -------------------------------------------------------

\subsection{System Model}

The scheme operates between a sender (encryptor) and receiver (decryptor) over an untrusted public channel under a threat model including quantum-capable adversaries. The security objective is to ensure that the shared key material cannot be recovered from the public channel by either a classical or quantum adversary, and that the cipher image is computationally indistinguishable from uniformly random noise. Three layers are stacked:
\begin{enumerate}
	\item \textit{Key Establishment:} ML-KEM-768 produces a session-unique 32-byte shared secret $K$ without transmitting it \cite{b6}.
	\item \textit{Key Derivation:} HKDF-SHA-256 expands $K$ into 144-byte OKM partitioned into 18 independent chaos parameters.
	\item \textit{Encryption:} A confusion-diffusion-confusion-diffusion pipeline using the Lorenz Attractor, Logistic Map, and Tent Map converts the plain image to a cipher image.
\end{enumerate}
The high-level data flow is: $K \xrightarrow{\text{HKDF-SHA-256}} \mathrm{OKM}[0..143] \xrightarrow{\text{map}} \theta_{1..18} \xrightarrow{\text{chaos}} E$.

\subsection{Key Establishment: ML-KEM-768}

The sender executes \texttt{ML-KEM.KeyGen()} to produce $\mathit{ek}$ (1,184~B) and $\mathit{dk}$ (2,400~B) \cite{b5}. The encapsulation key $\mathit{ek}$ is published over the public channel. The receiver runs \texttt{ML-KEM.Encaps($\mathit{ek}$)} to obtain ciphertext $c$ (1,088~B) and shared secret $K$ (32~B). The sender recovers the identical $K$ via \texttt{ML-KEM.Decaps($\mathit{dk}, c$)}. Because $K$ is never transmitted and its derivation is secured by the MLWE hardness assumption, the scheme provides quantum-resistant key establishment \cite{b6}.

\subsection{Key Derivation: HKDF-SHA-256}

The shared secret $K$ serves as Input Keying Material (IKM) for HKDF-SHA-256. A random salt and context label are applied in the Extract stage to produce a 32-byte pseudorandom key; five Expand iterations yield 160~bytes, of which the first 144 constitute the OKM. Each consecutive 8-byte block is converted to a 64-bit unsigned integer $w_i$ (big-endian), then normalized and linearly mapped to the valid chaotic parameter range:
\begin{equation}
	u_i = \frac{w_i}{2^{64}-1},\quad \theta_i = \theta_{\min} + u_i\,(\theta_{\max}-\theta_{\min}).
	\label{eq:parammap}
\end{equation}
Table~\ref{tab:hkdf_params} lists the full OKM byte allocation and valid ranges for all 18 parameters. Ranges are chosen to keep all maps within their established chaotic regimes \cite{b25,b28}.

\begin{table}[htbp]
	\caption{HKDF OKM Allocation and Chaotic Parameter Mapping}
	\label{tab:hkdf_params}
	\begin{center}
		\begin{tabular}{|c|l|l|l|}
			\hline
			\textbf{OKM Bytes} & \textbf{Parameter} & \textbf{Stage} & \textbf{Range} \\
			\hline
			0--7   & Lorenz $x_0$ & C1 & $[0.01,1.00)$ \\
			8--15  & Lorenz $y_0$ & C1 & $[0.01,1.00)$ \\
			16--23 & Lorenz $z_0$ & C1 & $[0.01,1.00)$ \\
			24--31 & Lorenz $\sigma$ & C1 & $[8.00,12.00)$ \\
			32--39 & Lorenz $\rho$ & C1 & $[24.00,35.00)$ \\
			40--47 & Lorenz $\beta$ & C1 & $[2.00,3.00)$ \\
			48--55 & Lorenz $dt$ & C1 & $[0.005,0.015)$ \\
			56--63 & Logistic $x_0$ & D1 & $[0.10,0.90)$ \\
			64--71 & Logistic $r$ & D1 & $[3.90,3.999)$ \\
			72--79 & Lorenz $x_0$ & C2 & $[0.01,1.00)$ \\
			80--87 & Lorenz $y_0$ & C2 & $[0.01,1.00)$ \\
			88--95 & Lorenz $z_0$ & C2 & $[0.01,1.00)$ \\
			96--103 & Lorenz $\sigma$ & C2 & $[8.00,12.00)$ \\
			104--111 & Lorenz $\rho$ & C2 & $[24.00,35.00)$ \\
			112--119 & Lorenz $\beta$ & C2 & $[2.00,3.00)$ \\
			120--127 & Lorenz $dt$ & C2 & $[0.005,0.015)$ \\
			128--135 & Tent $x_0$ & D2 & $[0.10,0.90)$ \\
			136--143 & Tent $\mu$ & D2 & $[1.90,1.999)$ \\
			\hline
		\end{tabular}
	\end{center}
\end{table}

\subsection{Encryption Pipeline}

The plain image $P$ of size $H\!\times\!W$ is flattened to a 1-D array $\mathbf{P}=[p_0,\ldots,p_{N-1}]$, $N=H\!\times\!W$. Four stages are then applied sequentially.

\textbf{Stage~1 --- 1st Confusion (Lorenz Attractor).}
The Lorenz system is integrated via Euler for $N$ steps from HKDF-derived $(x_0,y_0,z_0,\sigma,\rho,\beta,dt)$:
\begin{align}
	x_{n+1} &= x_n + \sigma(y_n - x_n)\,dt, \notag \\
	y_{n+1} &= y_n + [x_n(\rho - z_n) - y_n]\,dt, \label{eq:lorenz_euler} \\
	z_{n+1} &= z_n + (x_n y_n - \beta z_n)\,dt. \notag
\end{align}
Sequences $\{x_n\}$ and $\{y_n\}$ are argsorted to produce row- and column-permutation indices $(\pi_x, \pi_y)$, yielding the 1st-confused array $\mathbf{C}_1$.

\textbf{Stage~2 --- 1st Diffusion (Logistic Map).}
The Logistic Map is iterated $N$ times from HKDF-derived $(x_0, r)$ where $r \in [3.90, 3.999)$:
\begin{equation}
	x_{n+1} = r\,x_n(1-x_n).
	\label{eq:logistic_enc}
\end{equation}
Each floating-point value is converted to an 8-bit keystream byte:
\begin{equation}
	k_n = \lfloor\mathrm{frac}(|x_n|)\times256\rfloor\bmod 256.
	\label{eq:logistic_byte}
\end{equation}
The keystream $\mathbf{k}_1$ is XOR-combined with $\mathbf{C}_1$ to produce $\mathbf{D}_1 = \mathbf{C}_1 \oplus \mathbf{k}_1$.

\textbf{Stage~3 --- 2nd Confusion (Lorenz Attractor).}
A second Lorenz trajectory is integrated from a distinct non-overlapping OKM segment (bytes 72--127), producing a second independent permutation $(\pi_x^{(2)}, \pi_y^{(2)})$ with overwhelming probability \cite{b8,b19}. The 2nd confusion produces $\mathbf{C}_2$ from $\mathbf{D}_1$.

\textbf{Stage~4 --- 2nd Diffusion (Tent Map).}
The Tent Map is iterated $N$ times from HKDF-derived $(x_0, \mu)$ where $\mu \in [1.90, 1.999)$:
\begin{equation}
	x_{n+1} = \begin{cases} x_n/\mu, & x_n < \mu/2 \\ (1-x_n)/(1-\mu/2), & x_n \geq \mu/2. \end{cases}
	\label{eq:tent_enc}
\end{equation}
Its near-uniform output produces keystream $\mathbf{k}_2$ and final cipher array $\mathbf{E}=\mathbf{C}_2\oplus\mathbf{k}_2$, reshaped to $H\!\times\!W$. Because the Tent Map output is nearly uniformly distributed \cite{b25}, the resulting keystream produces a flat cipher-image histogram.

\subsection{Decryption}

Decryption is the exact inverse of encryption. The receiver regenerates all 18 chaos parameters from the same $K$ via HKDF, then applies: (1)~inverse 2nd diffusion ($\mathbf{C}_2 = \mathbf{E}\oplus\mathbf{k}_2$), (2)~inverse 2nd confusion ($(\pi_x^{(2)})^{-1}$, $(\pi_y^{(2)})^{-1}$), (3)~inverse 1st diffusion ($\mathbf{C}_1 = \mathbf{D}_1\oplus\mathbf{k}_1$), and (4)~inverse 1st confusion ($\pi_x^{-1}$, $\pi_y^{-1}$), recovering $\mathbf{P}$ exactly. A single-bit deviation in any chaos parameter yields a completely uncorrelated output, confirming chaotic sensitivity \cite{b28,b30}.

% -------------------------------------------------------
\section{Experimental Results}
% -------------------------------------------------------

\subsection{Setup and Correctness}

Nine grayscale PNG images in three resolution classes were used: $512\!\times\!512$ (Babon, Camar, Kubah), $900\!\times\!900$ (Gedung, Poster, Abstrak), and $1080\!\times\!1080$ (Bunga, Daun, Matahari), covering dark-dominant, bright-dominant, and balanced intensity distributions. The scheme was implemented in Python using \texttt{liboqs} bindings for ML-KEM-768 and the \texttt{cryptography} library for HKDF-SHA-256. All nine images were decrypted with bit-exact accuracy; a controlled single-bit fault injection into the Logistic Map initial condition $x_0$ rendered the recovered image completely unrecognizable, confirming that the scheme inherits the sensitive-dependence-on-initial-conditions property of the underlying chaotic maps \cite{b28,b30}.

\subsection{Visual Results}

Figs.~\ref{fig:babon}--\ref{fig:matahari} show the plain image, 1st confusion output, 1st diffusion output, 2nd confusion output, and 2nd diffusion (cipher image) for three representative images across the three resolution classes. Spatial structure is progressively eliminated; the final cipher images are visually indistinguishable from random noise.

% --- Figure: Babon (512x512) ---
\begin{figure*}[t]
	\centering
	\begin{minipage}[b]{0.18\linewidth}
		\centering
		\includegraphics[width=\linewidth]{babon.png}\\[2pt]
		{\footnotesize Original}
	\end{minipage}\hfill
	\begin{minipage}[b]{0.18\linewidth}
		\centering
		\includegraphics[width=\linewidth]{babon_confusion1.png}\\[2pt]
		{\footnotesize Confusion 1}
	\end{minipage}\hfill
	\begin{minipage}[b]{0.18\linewidth}
		\centering
		\includegraphics[width=\linewidth]{babon_diffusion1.png}\\[2pt]
		{\footnotesize Diffusion 1}
	\end{minipage}\hfill
	\begin{minipage}[b]{0.18\linewidth}
		\centering
		\includegraphics[width=\linewidth]{babon_confusion2.png}\\[2pt]
		{\footnotesize Confusion 2}
	\end{minipage}\hfill
	\begin{minipage}[b]{0.18\linewidth}
		\centering
		\includegraphics[width=\linewidth]{babon_diffusion2.png}\\[2pt]
		{\footnotesize Diffusion 2}
	\end{minipage}
	\caption{Encryption stages for Babon ($512\times512$). Left to right: original, after 1st confusion, after 1st diffusion, after 2nd confusion, cipher image (2nd diffusion).}
	\label{fig:babon}
\end{figure*}

% --- Figure: Poster (900x900) ---
\begin{figure*}[t]
	\centering
	\begin{minipage}[b]{0.18\linewidth}
		\centering
		\includegraphics[width=\linewidth]{poster.png}\\[2pt]
		{\footnotesize Original}
	\end{minipage}\hfill
	\begin{minipage}[b]{0.18\linewidth}
		\centering
		\includegraphics[width=\linewidth]{poster_confusion1.png}\\[2pt]
		{\footnotesize Confusion 1}
	\end{minipage}\hfill
	\begin{minipage}[b]{0.18\linewidth}
		\centering
		\includegraphics[width=\linewidth]{poster_diffusion1.png}\\[2pt]
		{\footnotesize Diffusion 1}
	\end{minipage}\hfill
	\begin{minipage}[b]{0.18\linewidth}
		\centering
		\includegraphics[width=\linewidth]{poster_confusion2.png}\\[2pt]
		{\footnotesize Confusion 2}
	\end{minipage}\hfill
	\begin{minipage}[b]{0.18\linewidth}
		\centering
		\includegraphics[width=\linewidth]{poster_diffusion2.png}\\[2pt]
		{\footnotesize Diffusion 2}
	\end{minipage}
	\caption{Encryption stages for Poster ($900\times900$).}
	\label{fig:poster}
\end{figure*}

% --- Figure: Matahari (1080x1080) ---
\begin{figure*}[t]
	\centering
	\begin{minipage}[b]{0.18\linewidth}
		\centering
		\includegraphics[width=\linewidth]{matahari.png}\\[2pt]
		{\footnotesize Original}
	\end{minipage}\hfill
	\begin{minipage}[b]{0.18\linewidth}
		\centering
		\includegraphics[width=\linewidth]{matahari_confusion1.png}\\[2pt]
		{\footnotesize Confusion 1}
	\end{minipage}\hfill
	\begin{minipage}[b]{0.18\linewidth}
		\centering
		\includegraphics[width=\linewidth]{matahari_diffusion1.png}\\[2pt]
		{\footnotesize Diffusion 1}
	\end{minipage}\hfill
	\begin{minipage}[b]{0.18\linewidth}
		\centering
		\includegraphics[width=\linewidth]{matahari_confusion2.png}\\[2pt]
		{\footnotesize Confusion 2}
	\end{minipage}\hfill
	\begin{minipage}[b]{0.18\linewidth}
		\centering
		\includegraphics[width=\linewidth]{matahari_diffusion2.png}\\[2pt]
		{\footnotesize Diffusion 2}
	\end{minipage}
	\caption{Encryption stages for Matahari ($1080\times1080$).}
	\label{fig:matahari}
\end{figure*}

\subsection{Histogram Analysis}

Histogram uniformity is a standard first-order test of encryption quality \cite{b15}. The standard deviation $\sigma_H$ of per-intensity-level frequency over all 256 gray levels was computed for both plain and cipher images; lower $\sigma_H$ indicates greater uniformity. Plain images exhibit highly non-uniform histograms ($\sigma_H = 2{,}491$--$6{,}784$), reflecting the spatial redundancy of natural images \cite{b15}. After encryption, all cipher images achieve dramatically reduced standard deviations ($\sigma_H = 36$--$72$), confirming near-uniform pixel intensity distribution consistently across all resolution classes (Table~\ref{tab:histogram}). These results confirm that the dual confusion--dual diffusion pipeline eliminates exploitable intensity-level patterns from the plaintext.

\begin{table}[htbp]
	\caption{Histogram Standard Deviation: Plain vs.\ Cipher Image}
	\label{tab:histogram}
	\begin{center}
		\begin{tabular}{|l|c|c|c|}
			\hline
			\textbf{Image} & \textbf{Size} & \textbf{Plain $\sigma_H$} & \textbf{Cipher $\sigma_H$} \\
			\hline
			Babon    & $512^2$   & 3506.90 & 43.51 \\
			Camar    & $512^2$   & 2491.33 & 36.23 \\
			Kubah    & $512^2$   & 3245.15 & 49.89 \\
			Gedung   & $900^2$   & 4088.65 & 56.65 \\
			Poster   & $900^2$   & 6783.65 & 57.08 \\
			Abstrak  & $900^2$   & 3322.11 & 54.79 \\
			Bunga    & $1080^2$  & 4781.00 & 71.59 \\
			Daun     & $1080^2$  & 3344.75 & 63.96 \\
			Matahari & $1080^2$  & 5711.30 & 67.64 \\
			\hline
		\end{tabular}
	\end{center}
\end{table}

\subsection{Pixel Correlation, Entropy, and PSNR}

Three additional metrics were computed. The horizontal adjacent-pixel correlation coefficient is:
\begin{equation}
	r = \frac{\sum_i (p_i - \bar{p})(p_{i+1} - \bar{p})}
	{\sqrt{\sum_i (p_i - \bar{p})^2 \cdot \sum_i (p_{i+1} - \bar{p})^2}},
	\label{eq:corr}
\end{equation}
with $|r| \approx 0$ indicating negligible linear dependence between adjacent pixels \cite{b14}. Information entropy is:
\begin{equation}
	H = -\sum_{v=0}^{255} p(v)\log_2 p(v),
	\label{eq:entropy}
\end{equation}
with theoretical maximum $H_\mathrm{max} = 8$ bits/symbol for a uniformly distributed 8-bit image. PSNR between plain and cipher image is:
\begin{equation}
	\mathrm{PSNR} = 10\log_{10}\!\left(\frac{255^2}{\mathrm{MSE}}\right)[\mathrm{dB}],
	\label{eq:psnr}
\end{equation}
where low PSNR confirms high visual dissimilarity \cite{b12}. Full results are in Table~\ref{tab:metrics}.

\begin{table}[htbp]
	\caption{Statistical Metrics of Cipher Images}
	\label{tab:metrics}
	\begin{center}
		\begin{tabular}{|l|c|c|c|}
			\hline
			\textbf{Image} & \textbf{Corr.} & \textbf{Entropy (bits)} & \textbf{PSNR (dB)} \\
			\hline
			Babon    & $-0.003527$ & 7.9987 & 8.9972 \\
			Camar    & $-0.001506$ & 7.9991 & 7.9632 \\
			Kubah    &  $0.000583$ & 7.9983 & 8.7849 \\
			Gedung   & $-0.000630$ & 7.9998 & 7.1152 \\
			Poster   &  $0.000438$ & 7.9998 & 6.0932 \\
			Abstrak  & $-0.000475$ & 7.9998 & 7.8507 \\
			Bunga    & $-0.001749$ & 7.9998 & 7.1025 \\
			Daun     &  $0.000061$ & 7.9999 & 7.8575 \\
			Matahari & $-0.000033$ & 7.9998 & 8.3739 \\
			\hline
			\textbf{Mean} & $\mathbf{-0.000715}$ & \textbf{7.9994} & \textbf{7.79} \\
			\hline
		\end{tabular}
	\end{center}
\end{table}

\textit{Correlation.} All cipher images achieve $|r| < 3.6\!\times\!10^{-3}$, with mean absolute value $\approx 7\!\times\!10^{-4}$. Both positive and negative values near zero confirm that the two-stage Lorenz confusion has randomized pixel ordering in both row and column directions \cite{b29}, consistent with results reported for multi-map confusion schemes \cite{b22,b26}.

\textit{Entropy.} All nine cipher images achieve $H > 7.998$ bits (mean 7.9994, peak 7.9999). The mean deviates from the theoretical maximum by only 0.0006 bits, confirming that dual diffusion distributes pixel intensities uniformly across all 256 levels. The Tent Map's near-uniform output distribution \cite{b25} prevents residual intensity concentration that would otherwise arise from the Logistic Map's non-uniform marginal.

\textit{PSNR.} Values range from 6.09~dB (Poster) to 8.99~dB (Babon), well below the 30~dB visual similarity threshold \cite{b12}. Variation across images reflects differences in plain-image pixel variance, not any structural weakness of the scheme.

% -------------------------------------------------------
\section{Discussion}
% -------------------------------------------------------

The experimental results demonstrate three key properties of the proposed scheme.

\textbf{Post-quantum security of key establishment.} By employing ML-KEM-768 (NIST Category~3), the shared secret used to seed all chaos parameters is protected against both classical and quantum adversaries---a property absent in all closely related prior works.

\textbf{Secure and independent parameter derivation.} The use of HKDF with distinct \texttt{info} labels for each chaos component ensures that the four sets of chaos parameters are cryptographically independent. This prevents cross-correlation between confusion and diffusion keys and eliminates key reuse across sessions.

\textbf{High-quality cipher images.} The dual confusion--dual diffusion architecture produces cipher images with near-zero pixel correlation, near-maximum entropy, and low PSNR across all images and resolutions. The Tent Map's more uniform output distribution \cite{b25} makes it the appropriate final-stage choice, while the Lorenz Attractor's 3-D structure provides a richer permutation space than any 1-D map. The current implementation is limited to grayscale images; extension to RGB color images and addition of a MAC-based authentication layer are identified as immediate future work.

% -------------------------------------------------------
\section{Conclusion}
% -------------------------------------------------------

This paper proposed and evaluated a hybrid Chaos-PQC image encryption scheme integrating ML-KEM-768 for quantum-resistant key establishment, HKDF-SHA-256 for cryptographically secure parameter derivation, and a four-stage confusion-diffusion-confusion-diffusion pipeline using the Lorenz Attractor, Logistic Map, and Tent Map. Experiments on nine grayscale images demonstrate that all cipher images achieve entropy ${>}\,7.998$ bits, pixel correlation ${<}\,|0.004|$, and PSNR ${<}\,9$~dB, confirming strong encryption quality and resistance to statistical analysis. The post-quantum security layer provided by ML-KEM-768 distinguishes this scheme from existing chaos-based approaches and makes it suitable for deployment in environments where long-term quantum-resistant security is required.

% -------------------------------------------------------
\section*{Acknowledgment}
% -------------------------------------------------------

The authors thank Yeni Farida, S.Stat., M.Si., for her guidance throughout the research, and acknowledge Politeknik Siber dan Sandi Negara (Poltek SSN) for the institutional support.

% -------------------------------------------------------
\begin{thebibliography}{00}

	\bibitem{b1}
	Z. Whittaker, ``A billion medical images are exposed online, as doctors ignore warnings,'' \textit{TechCrunch}, Jan. 2020.

	\bibitem{b2}
	R. K. Nistanto, ``Ini Hasil Pertemuan Kominfo dengan BPJS Kesehatan Terkait Kebocoran Data,'' \textit{KOMPAS.com}, May 2021.

	\bibitem{b3}
	R. Ramadhan, ``102 Juta Data Penduduk Indonesia di Kemensos Bocor, Dijual di Forum Hacker,'' \textit{Kumparan}, Sep. 2022.

	\bibitem{b4}
	Republik Indonesia, \textit{Undang-Undang Nomor 27 Tahun 2022 tentang Pelindungan Data Pribadi}, 2022.

	\bibitem{b5}
	C. Boutin, ``NIST releases first 3 finalized post-quantum encryption standards,'' NIST, Aug. 2024.

	\bibitem{b6}
	G. M. Raimondo et al., \textit{FIPS 203: Module-Lattice-Based Key-Encapsulation Mechanism Standard}, NIST, 2024.

	\bibitem{b7}
	E. Barker et al., \textit{NIST SP 800-56C Rev.~2: Recommendation for Key-Derivation Methods in Key-Establishment Schemes}, NIST, 2020.

	\bibitem{b8}
	H. Krawczyk and P. Eronen, ``HMAC-based Extract-and-Expand Key Derivation Function (HKDF),'' RFC~5869, IETF, 2010.

	\bibitem{b9}
	M. Kalra, S. Katyal, and R. Singh, ``A Tent map and Logistic map based approach for chaos-based image encryption,'' \textit{Lect. Notes Netw. Syst.}, 2019.

	\bibitem{b10}
	J. Giesl et al., ``Chaos-Based Bit Planes Image Encryption,'' Tomas Bata University.

	\bibitem{b11}
	S. Chen et al., ``Logistic map: Stability and entrance to chaos,'' \textit{J. Phys.: Conf. Ser.}, 2021.

	\bibitem{b12}
	S. T. Kamal et al., ``A new image encryption algorithm for grey and color medical images,'' \textit{IEEE Access}, 2021.

	\bibitem{b13}
	V. Vu\v{c}kovi\'{c}, ``Image and its matrix, matrix and its image,'' \textit{Review of the National Center for Digitization}, vol.~12, pp.~17--31, 2008.

	\bibitem{b14}
	S. Khaitan, S. Sagar, and R. Agarwal, ``Chaos cryptosystem with optimal key selection for image encryption,'' \textit{Multimedia Tools Appl.}, 2022.

	\bibitem{b15}
	P. Sarosh, S. A. Parah, and G. M. Bhat, ``An efficient image encryption scheme for healthcare applications,'' \textit{Multimedia Tools and Applications}, vol.~81, pp.~7253--7270, 2022, doi:~10.1007/s11042-021-11812-0.

	\bibitem{b17}
	M. Kaur and V. Kumar, ``A comprehensive review on image encryption techniques,'' \textit{Archives of Computational Methods in Engineering}, vol.~27, pp.~15--43, 2020.

	\bibitem{b18}
	J.-P. Aumasson, ``The impact of quantum computing on cryptography,'' \textit{Computer Fraud \& Security}, vol.~2017, no.~6, pp.~8--11, 2017.

	\bibitem{b19}
	H. Krawczyk, ``Cryptographic extraction and key derivation: The HKDF scheme,'' in \textit{Proc. 30th Annu. Cryptology Conf. (CRYPTO 2010)}, LNCS, vol.~6223, Springer, 2010, pp.~631--648.

	\bibitem{b20}
	Z. He et al., ``Design and analysis of a secure image encryption algorithm using proposed non-linear RN chaotic system and ECC/HKDF key derivation with authentication support,'' \textit{Scientific Reports}, vol.~15, art.~no.~39951, 2025.

	\bibitem{b21}
	J. M. McGinthy and A. J. Michaels, ``Further analysis of PRNG-based key derivation functions,'' \textit{IEEE Access}, vol.~7, pp.~95978--95986, 2019.

	\bibitem{b22}
	A. Qayyum et al., ``Chaos-based confusion and diffusion of image pixels using dynamic substitution,'' \textit{IEEE Access}, vol.~8, pp.~140876--140893, 2020.

	\bibitem{b23}
	K. H. Shakib, M. A. Hafeez, and A. Munir, ``AmphiKey: A dual-mode secure authenticated key encapsulation protocol for smart grid,'' arXiv:2509.01701, Sep. 2025.

	\bibitem{b24}
	Z. Hua and Y. Zhou, ``Dynamic parameter-control chaotic system,'' \textit{IEEE Trans. Cybern.}, vol.~46, no.~12, pp.~3330--3341, 2016.

	\bibitem{b25}
	M. Lawnik, ``Combined logistic and tent map,'' \textit{J. Phys.: Conf. Ser.}, vol.~1141, art.~no.~012132, 2018.

	\bibitem{b26}
	S. P. Praveen et al., ``A novel dual confusion and diffusion approach for grey image encryption using multiple chaotic maps,'' \textit{Int. J. Advanced Computer Science and Applications}, vol.~14, no.~8, 2023.

	\bibitem{b27}
	NIST, ``Security (evaluation criteria),'' Post-Quantum Cryptography Standardization. [Online]. Available: \texttt{csrc.nist.gov/Projects/Post-Quantum-Cryptography}.

	\bibitem{b28}
	E. N. Lorenz, ``Deterministic nonperiodic flow,'' \textit{J. Atmos. Sci.}, vol.~20, no.~3, pp.~130--141, 1963.

	\bibitem{b29}
	S. P. Praveen et al., ``A novel dual confusion and diffusion approach for grey image encryption using multiple chaotic maps,'' \textit{Int. J. Advanced Computer Science and Applications}, vol.~14, no.~8, pp.~836--843, 2023.

	\bibitem{b30}
	S. H. Strogatz, \textit{Nonlinear Dynamics and Chaos}, 2nd ed. Boulder, CO: Westview Press, 2015.

\end{thebibliography}

\end{document}
